\uuid{9Kaq}
\exo7id{1086}
\titre{exo7 1086}
\auteur{legall}
\organisation{exo7}
\datecreate{1998-09-01}
\isIndication{false}
\isCorrection{false}
\chapitre{Matrice}
\sousChapitre{Matrice et application linéaire}
\cpge{mpsi-mp2i.espaces-vectoriels.c, pcsi.espaces-vectoriels.c, ptsi.espaces-vectoriels.c, mpsi-mp2i.matrices-applications-lineaires.a, pcsi.matrices-applications-lineaires.a, ptsi.matrices-applications-lineaires.a, bcpst1.espace-vectoriel-kn.b, bcpst1.applications-lineaires-matrices}
\module{Algèbre}
\niveau{L1}
\difficulte{}

\contenu{
\texte{
Soient  $E$  un espace vectoriel de dimension  $n$, $f$  une application
lin\' eaire de  $E$  dans lui-m\^eme et  $x$  un \' el\' ement de  $E$  tel que la
famille  $f(x), ...,f^n(x)$  soit libre.
}
\begin{enumerate}
    \item \question{Montrer que la famille  $x,f(x),\ldots,f^{n-1}(x)$  est
une base de  $E$. D\' eduiser-en que  $f$  est bijective.}
    \item \question{On suppose maintenant que  $f^n(x)=x$. D\' eterminer la
matrice de  $f$  dans la
base   $x,f(x),\ldots,f^{n-1}(x)$.}
\end{enumerate}
}
